% !TEX program = xelatex
% ============================================================
%  Chapter 1 Deep Generative Modeling — English Main Version
% ============================================================

\documentclass[11pt,a4paper]{article}

% ------------------------------------------------------------
%  Typography and page layout
% ------------------------------------------------------------
\IfFileExists{geometry.sty}{%
  \usepackage[margin=2.5cm]{geometry}
}{%
  \PackageWarning{geometry}{geometry not found, fallback to default margins.}%
}
\usepackage{ctex}
\usepackage{microtype}
\linespread{1.08}

% ------------------------------------------------------------
%  Utility packages
% ------------------------------------------------------------
\usepackage{amsmath,amssymb,amsthm}
\usepackage{mathtools}
\usepackage{graphicx}
\usepackage{booktabs}
\usepackage{enumitem}
\usepackage[dvipsnames]{xcolor}
\usepackage[most]{tcolorbox}
\usepackage{titlesec}
\setlength{\parskip}{0.35em}
\setlength{\parindent}{2em}
\setlist[itemize]{itemsep=0.2em, topsep=0.25em}
\setlist[enumerate]{itemsep=0.2em, topsep=0.25em}
\newcommand{\bx}{\mathbf{x}}

% ------------------------------------------------------------
%  Colors and highlighting
% ------------------------------------------------------------
\definecolor{mydarkblue}{rgb}{0.0,0.08,0.45}
\definecolor{cardbg}    {RGB}{245,245,247}
\definecolor{highlight} {HTML}{FCD66F}
\definecolor{warning}   {HTML}{E57373}
\definecolor{info}      {HTML}{4DB6AC}
\newcommand{\tlight}[1]{\textcolor{warning}{#1}}
\newcommand{\tinfo} [1]{\textcolor{info}{#1}}
\newcommand{\cbox}[2]{\colorbox{#1}{\textcolor{black}{#2}}}
\newcommand{\bglight}[1]{\cbox{highlight}{#1}}
\newcommand{\bgwarn} [1]{\cbox{warning}{#1}}
\newcommand{\bginfo} [1]{\cbox{info}{#1}}

% ------------------------------------------------------------
%  Hyperlinks
% ------------------------------------------------------------
\usepackage{hyperref}
\hypersetup{
  colorlinks = true,
  linkcolor  = mydarkblue,
  citecolor  = mydarkblue,
  urlcolor   = mydarkblue
}

% ------------------------------------------------------------
%  Section styles
% ------------------------------------------------------------
\titleformat{\section}
  {\large\bfseries\color{mydarkblue}}{\thesection}{0.5em}{}[\titlerule]
\titlespacing*{\section}{0pt}{1.5ex}{1ex}
\titleformat{\subsection}
  {\normalsize\bfseries}{\thesubsection}{0.5em}{}
\titlespacing*{\subsection}{0pt}{1ex}{0.5ex}

% ------------------------------------------------------------
%  Theorem environments (tcolorbox)
% ------------------------------------------------------------
\tcbset{
  mytheo/.style={
    enhanced, breakable,
    colback=#1!5!white, colframe=#1!70!black,
    coltitle=white, fonttitle=\bfseries\small,
    boxrule=0.6pt, arc=3pt,
    left=5pt, right=5pt, top=10pt, bottom=5pt,
    attach boxed title to top left={yshift*=-\tcboxedtitleheight/2, xshift=2mm},
    boxed title style={
      colback=#1!80!black, colframe=#1!80!black,
      boxrule=0pt, arc=2pt
    }
  }
}
\newtcbtheorem[auto counter]             {dgmdef}    {Definition}{mytheo=NavyBlue}{dgmdef}
\newtcbtheorem[use counter from=dgmdef]  {dgmthm}    {Theorem}{mytheo=WildStrawberry}{dgmthm}
\newtcbtheorem[use counter from=dgmdef]  {dgmremark} {Remark}{mytheo=Periwinkle}{dgmrem}

\tcolorboxenvironment{proof}{
  blanker, breakable, left=4pt, right=4pt, top=3pt, bottom=3pt,
  borderline west={2pt}{0pt}{gray!50}
}

% ------------------------------------------------------------
%  Header + horizontal navigation
% ------------------------------------------------------------
\makeatletter
\newcommand{\makeDGMTitleEN}{%
  \begin{tcolorbox}[
    colback=mydarkblue!5, colframe=mydarkblue!10,
    arc=3pt, boxrule=0.5pt,
    left=5pt, right=5pt, top=5pt, bottom=5pt
  ]
    {\large\bfseries\@title}\hfill{\small\@author}\\[-2pt]
    {\tiny\color{gray}\@date\hfill Chapter 1 · Diffusion and Score Matching}
  \end{tcolorbox}
  \vspace{-0.6em}
}
\makeatother

\newtcolorbox{tabNavBoxSmall}[1][]{
  enhanced,
  colframe=white, colback=gray!5, coltext=black!80,
  fontupper=\sffamily\small\bfseries,
  arc=1.5mm, boxrule=0pt, bottomrule=2pt,
  left=2mm, right=2mm, top=1.5mm, bottom=1.5mm,
  #1
}

\newcommand{\horizontalTOCDGMEN}{%
  \begin{center}
  \begin{tcbraster}[
    raster columns=3, raster width=\linewidth,
    raster column skip=1mm, raster equal height
  ]
    \begin{tabNavBoxSmall}
      \hyperref[sec:setup]{\ Setup and Goal}
    \end{tabNavBoxSmall}
    \begin{tabNavBoxSmall}
      \hyperref[sec:divergence]{\ Divergences and Training}
    \end{tabNavBoxSmall}
    \begin{tabNavBoxSmall}
      \hyperref[sec:taxonomy]{\ Taxonomy and Summary}
    \end{tabNavBoxSmall}
  \end{tcbraster}
  \vspace{-2.5pt}%
  \noindent\textcolor{gray!30}{\rule{\linewidth}{1pt}}
  \end{center}
  \vspace{1.2em}
}

\title{Diffusion and Score Matching \\ \large Chapter 1: Deep Generative Modeling Basics (English-main)}
\author{}
\date{\today}

\begin{document}

\makeDGMTitleEN
\horizontalTOCDGMEN

\begin{center}
  \includegraphics[width=\linewidth]{\detokenize{ChatGPT Image 2026年6月5日 20_07_56.png}}
\end{center}
\vspace{0.5em}

\section{Deep Generative Modeling}\label{sec:setup}
\subsection{Problem Setup}
\textbf{Input:}
A collection of real-world samples drawn from an unknown data distribution $P_{\mathrm{data}}(\bx)$.

\textbf{Output:}
A neural-network parameterized model distribution $P_\phi(\bx)$.

\textbf{Goal:}
\begin{enumerate}[label=(\arabic*)]
  \item Generate samples that resemble real data;
  \item Enable controllable generation.
\end{enumerate}

\subsection{Mathematical Formulation}
Let the training set be
\[
\mathcal D=\{\bx^{(i)}\}_{i=1}^{N},\qquad \bx^{(i)}\overset{\mathrm{i.i.d.}}{\sim}P_{\mathrm{data}}.
\]
We seek parameters $\phi^\star$ such that
\[
P_{\phi^\star}(\bx)\approx P_{\mathrm{data}}(\bx).
\]
Training is framed as minimizing a discrepancy between distributions:
\[
\phi^\star=\arg\min_\phi D\!\big(P_{\mathrm{data}},P_\phi\big).
\]

\section{Divergences and Training Objectives}\label{sec:divergence}
\subsection{KL Divergence (Forward KL)}

\begin{dgmdef}{Kullback--Leibler Divergence}{kl-divergence}
\[
D_{\mathrm{KL}}(P_{\mathrm{data}}\parallel P_\phi)
\triangleq \int P_{\mathrm{data}}(\bx)\,
\log\frac{P_{\mathrm{data}}(\bx)}{P_\phi(\bx)}\,d\bx
=\mathbb{E}_{\bx\sim P_{\mathrm{data}}}
\left[\log\frac{P_{\mathrm{data}}(\bx)}{P_\phi(\bx)}\right].
\]
\end{dgmdef}

Expanding yields
\[
D_{\mathrm{KL}}(P_{\mathrm{data}}\parallel P_\phi)
=\mathbb{E}_{\bx\sim P_{\mathrm{data}}}
\big[\log P_{\mathrm{data}}(\bx)-\log P_\phi(\bx)\big]
=-\mathbb{E}_{\bx\sim P_{\mathrm{data}}}\big[\log P_\phi(\bx)\big]-H(P_{\mathrm{data}}),
\]
where
\[
H(P_{\mathrm{data}})
\triangleq -\mathbb{E}_{\bx\sim P_{\mathrm{data}}}\big[\log P_{\mathrm{data}}(\bx)\big]
\]
is the entropy of the data distribution and does not depend on $\phi$. Therefore:

\begin{dgmthm}{Forward KL $\Longleftrightarrow$ Maximum Likelihood Estimation}{kl-mle}
\[
\min_\phi D_{\mathrm{KL}}(P_{\mathrm{data}}\parallel P_\phi)
\Longleftrightarrow
\max_\phi \mathbb{E}_{\bx\sim P_{\mathrm{data}}}[\log P_\phi(\bx)].
\]
With finite data, this corresponds to the empirical objective
\[
\max_\phi \frac{1}{N}\sum_{i=1}^{N}\log P_\phi\!\big(\bx^{(i)}\big).
\]
\end{dgmthm}

\subsection{Fisher Divergence}

\begin{dgmdef}{Fisher Divergence}{fisher-divergence}
\[
D_F(P\parallel Q)
=\mathbb{E}_{\bx\sim P}
\left[\left\|\nabla_{\bx}\log P(\bx)-\nabla_{\bx}\log Q(\bx)\right\|^2\right].
\]
This divergence is always nonnegative: $D_F(P\parallel Q)\ge 0$. Under suitable regularity conditions,
\[
D_F(P\parallel Q)=0 \iff \nabla_{\bx}\log P(\bx)=\nabla_{\bx}\log Q(\bx)\quad (P\text{-a.s.}).
\]
\end{dgmdef}

In diffusion/score matching settings, one typically sets $P=P_{\mathrm{data}}$ and $Q=P_\phi$, so training matches score functions.

\subsection{$f$-Divergence Family}

\begin{dgmdef}{$f$-Divergence}{f-divergence}
For convex $f:\mathbb{R}_+\to\mathbb{R}$ with $f(1)=0$:
\[
D_f(P\parallel Q)=\int Q(\bx)\,f\!\left(\frac{P(\bx)}{Q(\bx)}\right)\,d\bx.
\]
\end{dgmdef}

Common members:
\begin{itemize}
  \item \textbf{Jensen--Shannon Divergence}
  \[
  f_{\mathrm{JS}}(u)=u\log u-(u+1)\log\frac{u+1}{2},
  \]
  \[
  D_{\mathrm{JS}}(P\parallel Q)
  =\frac12 D_{\mathrm{KL}}(P\parallel M)+\frac12 D_{\mathrm{KL}}(Q\parallel M),
  \quad M=\frac{P+Q}{2}.
  \]

  \item \textbf{Total Variation Distance}
  \[
  f_{\mathrm{TV}}(u)=\frac12|u-1|,
  \]
  \[
  D_{\mathrm{TV}}(P,Q)
  =\frac12\int\!\left|P(\bx)-Q(\bx)\right|d\bx
  =\sup_A\left|P(A)-Q(A)\right|.
  \]
\end{itemize}

\subsection{Challenges in Explicit Density Modeling}
A valid density $P_\phi(\bx)$ must satisfy:
\begin{enumerate}
  \item Non-negativity: $P_\phi(\bx)\ge 0$;
  \item Normalization: $\int P_\phi(\bx)\,d\bx=1$.
\end{enumerate}
A common construction is the Energy-Based Model (EBM), written in two steps:
\begin{enumerate}
  \item Define an energy network $E_\phi:\mathbb{R}^D\to\mathbb{R}$ and obtain an unnormalized density
  \[
  \widetilde P_\phi(\bx)=\exp(E_\phi(\bx));
  \]
  \item Normalize it with the partition function
\[
P_\phi(\bx)=\frac{\exp(E_\phi(\bx))}{Z(\phi)},
\qquad
Z(\phi)=\int\exp(E_\phi(\bx))\,d\bx.
\]
\end{enumerate}
In practice, the partition function $Z(\phi)$ is usually intractable in high dimensions.

\section{Representative Generative Models}\label{sec:taxonomy}
\begin{enumerate}

  \item \textbf{Energy-Based Models (EBM)}

  Define an energy function $E_\phi:\mathbb{R}^D\to\mathbb{R}$ via a neural network (using the \textbf{negative-energy} convention):
  \[
  P_\phi(\bx)
  \triangleq \frac{1}{Z(\phi)}\exp\!\big({-}E_\phi(\bx)\big),
  \qquad
  Z(\phi)=\int\exp\!\big({-}E_\phi(\bx)\big)\,d\bx.
  \]
  The partition function $Z(\phi)$ is generally intractable in high dimensions.

  \item \textbf{Autoregressive Models}

  Factorize the joint distribution by the chain rule:
  \[
  P(\bx)=\prod_{i=1}^{D}P_\phi(x_i\mid x_{<i}),
  \qquad
  x_{<i}\triangleq(x_1,x_2,\ldots,x_{i-1}).
  \]
  Each conditional $P_\phi(x_i\mid x_{<i})$ is parameterized by a neural network (e.g.\ a Transformer).\\
  Training objective: minimize the negative log-likelihood
  \[
  \min_\phi\;-\frac{1}{N}\sum_{n=1}^{N}\sum_{i=1}^{D}\log P_\phi\!\big(x_i^{(n)}\mid x_{<i}^{(n)}\big).
  \]

  \item \textbf{Variational Autoencoders (VAE)}

  Introduce a latent variable $z$ to capture hidden structure in data:
  \begin{itemize}
    \item Encoder: $q_\phi(z\mid \bx)$
    \item Decoder: $p_\theta(\bx\mid z)$
  \end{itemize}
  Training objective: maximize the Evidence Lower Bound (ELBO)
  \[
  \mathcal{L}_{\mathrm{ELBO}}(\theta,\phi;\bx)
  =\mathbb{E}_{z\sim q_\phi(z\mid \bx)}\!\big[\log p_\theta(\bx\mid z)\big]
  -D_{\mathrm{KL}}\!\big(q_\phi(z\mid \bx)\parallel p(z)\big)
  \;\le\;\log p_\theta(\bx).
  \]

  \item \textbf{Normalizing Flows}

  Learn a \textbf{bijective mapping} $f$ between data space $\bx$ and latent space $z$:
  \[
  \bx \;\xrightarrow{f}\; z \;\xrightarrow{f^{-1}}\; \bx'.
  \]
  By the change-of-variables formula:
  \[
  \log P_\phi(\bx)
  =\log p_Z\!\big(f(\bx)\big)
  +\log\left|\det\frac{\partial f(\bx)}{\partial \bx}\right|.
  \]
  With a simple prior $p_Z$ (e.g.\ standard Gaussian), log-likelihood is computed exactly.

  \item \textbf{Generative Adversarial Networks (GAN)}

  GAN trains two networks adversarially:
  \begin{itemize}
    \item \textbf{Discriminator} $D$: distinguishes real from generated samples
    \[
    \bx \;\longrightarrow\; D(\bx) \;\longrightarrow\; 0/1
    \]
    \item \textbf{Generator} $G$: maps prior noise to fake samples
    \[
    z\sim p_{\mathrm{prior}}(z) \;\longrightarrow\; G(z) \;\longrightarrow\; \bx'
    \]
  \end{itemize}
  Standard minimax objective:
  \[
  \min_{G}\max_{D}
  \ \mathbb{E}_{\bx\sim P_{\mathrm{data}}}[\log D(\bx)]
  +\mathbb{E}_{z\sim p_{\mathrm{prior}}(z)}[\log(1-D(G(z)))].
  \]
  \textbf{Key point:} GAN does \textbf{not} define a density $P(\bx)$ — it is a canonical example of implicit modeling.

\end{enumerate}

\subsection{Explicit vs.\ Implicit Modeling}

\begin{center}
\renewcommand{\arraystretch}{1.35}
\begin{tabular}{@{}lll@{}}
\toprule
 & \textbf{Explicit Modeling} & \textbf{Implicit Modeling} \\
\midrule
Density $P(\bx)$ & Explicitly defined (or approximated) & Not defined \\
Likelihood & Computable (exactly or approximately) & Cannot be directly evaluated \\
Training & Maximize likelihood / ELBO & Adversarial training \\
Sampling & Via model distribution & $\bx=G_\phi(z),\ z\sim p_\mathrm{prior}$ \\
Examples & VAE, Normalizing Flows, AR, EBM & GAN \\
\bottomrule
\end{tabular}
\end{center}

\section{Summary}
\begin{itemize}
  \item The core objective is to align $P_\phi$ with $P_{\mathrm{data}}$;
  \item Forward KL leads to MLE but can be computationally challenging in high dimensions;
  \item Different model classes balance tractability and expressiveness differently;
  \item Diffusion models can be viewed as score matching plus iterative denoising, to be covered next.
\end{itemize}

\end{document}
